\section{Extension to smooth one-parameter atom families}
\label{app:general-atom}

The diversity--collision mechanism is not specific to exponentials.  Let
$\mathcal H$ be a finite-dimensional real Hilbert space, let $\mathcal L$ be
a fixed linear nuisance subspace, and let $\theta\mapsto f_\theta\in\mathcal H$
be five times continuously differentiable.  For $C$ channels define
\[
\mathcal M_1=\{\ell+s f_k^T:\ell\in\mathbb R^C\otimes\mathcal L,
s\in\mathbb R^C,k\in K\}
\]
and the two-atom mean
\[
M_{\theta,\delta}=\ell+a f_\theta^T+b f_{\theta+\delta}^T.
\]
As before, put $s=a+b$, $S=\|s\|^2$,
$p=\langle b,s\rangle/S$, $b_\perp=b-ps$, $B=\|b_\perp\|^2$, and
$c=p(1-p)/2$.  Let
\[
P_2=P_{(\mathcal L+\operatorname{span}\{f_\theta\})^\perp},\qquad
P_3=P_{(\mathcal L+\operatorname{span}\{f_\theta,f'_\theta\})^\perp},
\]
and define $V_\theta=\|P_2f'_\theta\|^2$ and
$W_\theta=\|P_3f''_\theta\|^2$.

\begin{theorem}[smooth-atom diversity--collision law]
\label{thm:general-atom}
Suppose $\theta$ ranges over a compact subset $K_0\subset\operatorname{int}K$,
$S>0$, $p\in[\epsilon,1-\epsilon]$, $B/S$ lies in a fixed compact subset of
$[0,\infty)$, and:
\begin{enumerate}
\item $f$ and its first five derivatives are uniformly bounded on a
neighborhood of $K_0$;
\item $V_\theta$ and $W_\theta$ are uniformly bounded away from zero;
\item the relevant nuisance-plus-atom and tangent Gram matrices have constant
rank, with their nonzero singular values uniformly bounded away from zero;
\item for every fixed $\eta>0$ small enough,
\[
\inf_{\theta\in K_0}\inf_{|k-\theta|\ge\eta}
\|P_{(\mathcal L+\operatorname{span}\{f_k\})^\perp}f_\theta\|>0.
\]
\end{enumerate}
Then, uniformly over the stated sets as $\delta\to0$,
\[
\operatorname{dist}^2(M_{\theta,\delta},\mathcal M_1)
=BV_\theta\delta^2+Sc^2W_\theta\delta^4
+O(B|\delta|^3+S|\delta|^5).
\]
Consequently the distance-based phase diagram remains valid after replacing
the exponential derivatives by $f'_\theta,f''_\theta$. The plug-in
phase-matched score limit additionally requires the null-chart smoothness,
tangent nonsingularity, amplitude lower bound, remote-fit separation, and
local-array conditions of Lemma~\ref{lem:fitted-frame} and
Theorem~\ref{thm:score}.
\end{theorem}

\begin{proof}
For $P_k=P_{(\mathcal L+\operatorname{span}\{f_k\})^\perp}$ and
$g_\delta=f_{\theta+\delta}-f_\theta$, exact profiling of the channelwise
linear coefficients gives
\[
D^2=S\min_{k\in K}\left\{
\|P_k(f_\theta+pg_\delta)\|^2
+{B\over S}\|P_kg_\delta\|^2\right\}.
\]
The cross term is zero because $s^Tb_\perp=0$. Assumption 4 localizes every
global minimizer to $k=\theta+o(1)$.  At $(k,\delta)=(\theta,0)$ the second
derivative in $k$ is $2V_\theta$, so Assumptions 1--3 and the uniform implicit
function theorem give a unique minimizing branch
$k_*=\theta+p\delta+O(\delta^2+(B/S)\delta^2)$.

Let $r=B/S$ and denote the profiled objective by $H(r,\delta)$.  The envelope
identity and $g_\delta=\delta f'_\theta+O(\delta^2)$ yield
\[
\partial_rH(r,\delta)=V_\theta\delta^2+O(|\delta|^3).
\]
Integration from zero to $r$ preserves the factor $r$ and gives
$H(r,\delta)-H(0,\delta)=rV_\theta\delta^2+O(r|\delta|^3)$.
When $r=0$, the fitted displacement $p\delta$ removes the first derivative.
After optimizing its second-order correction and all linear nuisance
coefficients, the residual at order $\delta^2$ is exactly
$c\delta^2P_3f''_\theta$.  Hence
$H(0,\delta)=c^2W_\theta\delta^4+O(|\delta|^5)$ uniformly.  Multiplication by
$S$ proves the result.
\end{proof}

The exponential family takes $\mathcal L=\operatorname{span}\{\mathbf1\}$
and $f_\lambda(t)=e^{-\lambda t}$.  Other covered examples include Gaussian
scale atoms $e^{-\theta t^2}$ and known-carrier damped atoms after stacking real
and imaginary parts, provided all rank and global no-alias conditions above
hold. Thus the theorem identifies a generic
osculating-curve phenomenon rather than an isolated exponential identity.
